% !TEX root = ../Preparazione alle olimpiadi.tex

\chapter{Programmazione Dinamica}

La \textbf{programmazione dinamica} è una particolare tecnica di programmazione che consente di trovare la soluzione a un problema scomponendolo in sottoproblemi più semplici, spesso sovrapposti tra loro, le cui soluzioni vengono memorizzate e riutilizzate per evitare calcoli ripetuti.

Questa tecnica può essere applicata quando un problema presenta una \textit{struttura di sottoproblemi ottimali} e quando i sottoproblemi non sono indipendenti, ma condividono parti comuni.

In questo modo, la programmazione dinamica combina la correttezza della ricerca esaustiva con l’efficienza di un approccio basato sulla memorizzazione dei risultati intermedi.

Esistono due principali utilizzi di questa tecnica:
\begin{itemize}
	\item \textbf{Trovare una soluzione ottimale}: si vuole determinare una soluzione che sia massima o minima rispetto a un certo criterio (ad esempio il costo minimo o il guadagno massimo).
	\item \textbf{Contare il numero di soluzioni}: si vuole calcolare il numero totale di soluzioni possibili che soddisfano determinate condizioni.
\end{itemize}



\section{Programmazione dinamica ricorsiva}

\section{Programmazione dinamica iterativa}